The difference between the simulated and experimental \gls{rmse} values
can be attributed to factors not considered in the \gls{fopdt} model used for
the controller design. The main sources of discrepancy include sensor
measurement noise, transport delays, and unmodeled thermal dynamics of the
physical plant.

Another source of discrepancy arises from the system architecture itself, in
which the controller uses the measurement from Sensor 1 for feedback, while the
error is computed based on the temperature measured by Sensor 2. The spatial
separation between the sensors introduces a delay between the controlled
variable and the variable of interest, which was not fully captured by the model
employed in the design process.

The disturbance introduced at 10 minutes, resulting from the activation of a
fan, also contributed to the increase in the experimental error, since its
actual effect on the plant differs from the simplified representation adopted in
the simulation.

Despite these discrepancies, the close agreement between the simulated and
experimental \gls{rmse} values confirms that the identified model adequately
represents the plant dynamics and that the proposed control strategy maintains
good reference-tracking performance under experimental conditions as well.
